% !TEX root = 00_main_aaai_2027.tex
% Appendix: the ranked probability score (RPS).
% Provenance: every worked-example number is printed by rps_appendix_examples.py in this directory;
% the experiment-11 numbers come from experiments/archive/11_formulation4_train_only_ce_mixture_cv_start_10Nov2026_end_10Nov2026/results.md.
\section{The Ranked Probability Score (RPS)}\label{app:rps}

\paragraph{Setting.}
A human rating is an ordered category $r\in\{0,1,\dots,5\}$.
For an artifact, a judge outputs a forecast distribution $q=(q_0,\dots,q_5)$ over ratings.
Its cumulative distribution function (CDF) and survival function are
\begin{align}
Q_j &\stackrel{\rm def}{=} \Pr_q[R\le j]=\sum_{k=0}^{j} q_k,\\
S_j &\stackrel{\rm def}{=} \Pr_q[R> j]=1-Q_j,
\end{align}
for thresholds $j=0,\dots,4$. The last threshold is omitted because $Q_5=1$ for every forecast.

\paragraph{Definition.}
The ranked probability score (RPS)~\citep{epstein1969scoring} compares the judge's CDF with the step function of the observed rating:
\begin{align}
L_{\rm RPS}(q,r) &\stackrel{\rm def}{=} \frac{1}{5}\sum_{j=0}^{4}\bigl(Q_j-\mathbf{1}\{r\le j\}\bigr)^2 \label{eq:rps-cdf}\\
&= \frac{1}{5}\sum_{j=0}^{4}\bigl(S_j-\mathbf{1}\{r> j\}\bigr)^2. \label{eq:rps-sf}
\end{align}
We use it as a loss: lower is better.
Dividing by the five thresholds puts it in $[0,1]$.
It is $0$ only when $q$ puts all its mass on $r$, and $1$ when $q$ puts all its mass on the opposite end of the scale (for example $q_0=1$ when $r=5$).
Some references omit the factor $1/5$; their values are five times ours.

\paragraph{Empirical RPS.}
For $N$ artifacts $x_i$ with $n_i$ human ratings $r_{i1},\dots,r_{in_i}$ each, we average over ratings and then artifacts:
\begin{equation}
\widehat{\rm RPS}=\frac{1}{N}\sum_{i=1}^{N}\frac{1}{n_i}\sum_{m=1}^{n_i}L_{\rm RPS}(q_i,r_{im}).
\end{equation}
Writing $\hat P_{ij}$ for the empirical CDF of artifact $i$'s ratings, the inner average equals $\frac{1}{5}\sum_j[(Q_{ij}-\hat P_{ij})^2+\hat P_{ij}(1-\hat P_{ij})]$.
One rating per artifact already gives an unbiased estimate of the judge's expected score; no per-artifact histogram has to be estimated.

\paragraph{Intuition 1: five yes/no questions.}
At each threshold $j$ the judge answers ``is the rating at most $j$?'' with probability $Q_j$, and the human rating supplies the answer.
Each summand in Eq.~\ref{eq:rps-cdf} is the Brier score of one such binary forecast, so RPS is the average Brier score over the five threshold questions.
Equivalently (Eq.~\ref{eq:rps-sf}), the judge forecasts the survival function: the chance that the rating exceeds each threshold.

\paragraph{Intuition 2: distance between staircases.}
The judge's CDF is a staircase that rises from $0$ to $1$; the observed rating is a single step at $r$.
RPS sums the squared vertical gaps between the two at each threshold.
A near miss leaves gaps at few thresholds and a far miss at many.
Table~\ref{tab:rps-examples} shows this for one human rating $r=4$.
Cross-entropy (CE) is infinite (or set by a probability floor) for every forecast with $q_4=0$, and the Brier score over the six categories, $\sum_k(q_k-\mathbf{1}\{r=k\})^2$, is $2$ for every point mass away from $4$. Only RPS gives the near miss a better (lower) score than the far miss.

\begin{table}[h]
\centering\small
\begin{tabular}{lccc}
\toprule
Forecast for $r=4$ & RPS & Brier & CE \\
\midrule
all mass on 4 & 0.000 & 0.0 & 0 \\
all mass on 5 & 0.200 & 2.0 & $\infty$ \\
all mass on 0 & 0.800 & 2.0 & $\infty$ \\
half on 3, half on 5 & 0.100 & 1.5 & $\infty$ \\
\bottomrule
\end{tabular}
\caption{Scores of four forecasts against one human rating $r=4$ (lower is better).}
\label{tab:rps-examples}
\end{table}

\paragraph{Intuition 3: rating points.}
For integer ratings, RPS has an absolute-error form, the discrete case of the identity for the continuous ranked probability score~\citep{matheson1976scoring,gneiting2007strictly}:
\begin{equation}
5\,L_{\rm RPS}(q,r)=\E_q\lvert R-r\rvert-\tfrac{1}{2}\E_q\lvert R-R'\rvert,
\end{equation}
where $R$ and $R'$ are independent draws from $q$.
The first term is the expected distance, in rating points, between a rating sampled from the judge and the human rating.
The second term subtracts half the judge's own spread, so a judge is not penalized merely for spreading its probability.

\paragraph{Why RPS is proper.}
Let the human rating be drawn from $p$ with CDF $P$. Then
\begin{align}
\E_{R\sim p}\bigl[L_{\rm RPS}(q,R)\bigr]
&=\underbrace{\frac{1}{5}\sum_{j=0}^{4}(Q_j-P_j)^2}_{\text{excess: judge's error}}\nonumber\\
&\quad+\underbrace{\frac{1}{5}\sum_{j=0}^{4}P_j(1-P_j)}_{\text{rater disagreement}}.
\label{eq:rps-decomposition}
\end{align}
The second term does not depend on $q$: it is the irreducible disagreement among raters.
The first term is nonnegative and is zero only when $Q=P$, that is, $q=p$.
So RPS is strictly proper~\citep{gneiting2007strictly}: a judge minimizes its expected RPS only by reporting the full human rating distribution, including how much raters disagree.

\paragraph{Example: disagreement that correlation cannot see.}
Suppose raters split evenly between ratings 3 and 5 on an artifact.
Judge~A predicts the human distribution; Judge~B puts all its mass on 4.
Both have expected rating 4, so on every such artifact they produce the same expected rating and any correlation with human mean ratings scores them identically.
Their expected RPS differs: $0.10$ for Judge~A (all irreducible disagreement) and $0.20$ for Judge~B. A uniform forecast over $0,\dots,5$ scores $0.206$.

\paragraph{Why RPS suits judges of ordinal human ratings.}
\begin{itemize}
\item \textbf{Order matters.} Ratings are ordered; RPS charges more for larger misses. CE and the categorical Brier score treat ratings as unordered labels.
\item \textbf{It targets the whole distribution.} Being strictly proper, RPS is minimized only by the true conditional rating distribution. Correlation ignores calibration and spread, and mean squared error identifies only the conditional mean.
\item \textbf{It works with sparse ratings.} Each observed rating is scored directly, so artifacts with one to three ratings contribute without estimating a histogram.
\item \textbf{It is bounded.} RPS lies in $[0,1]$ and cannot be dominated by one rating that received near-zero probability. CE can: in our retrospective human-rating analysis (archived experiment~11), judge mixtures had higher held-out CE than a constant forecast of the training ratings ($5.43$ vs $1.85$ nats), driven by the $10^{-6}$ probability floor, while their held-out RPS was lower ($0.137$ vs $0.183$).
\item \textbf{It is interpretable.} The absolute-error form is in rating points, and Eq.~\ref{eq:rps-decomposition} separates the judge's error from the noise floor set by rater disagreement.
\end{itemize}

\paragraph{Limitations.}
RPS assumes equally spaced ratings: a miss from 0 to 1 costs the same as a miss from 4 to 5. It weights all thresholds equally. It is less sensitive than CE to the probability given to rare ratings, which is why we keep CE as a secondary metric. With two or three ratings per artifact, $\hat P_{ij}(1-\hat P_{ij})$ is a noisy estimate of rater disagreement, so comparisons between judges should use paired differences on the same artifacts.
